% GENERATED FILE. Edit canonical lessons, then run npm run generate:books.
% canonical-source-hash: fnv1a64:c9d35de6063a5291
% canonical-lessons: RU-C04-dumat, RU-C04-ponimat, RU-C04-chitat, RU-C04-pisat

\chapter{Verbs of the Mind and the Page}
\label{ch:russian-verbs-of-the-mind}
\hlchaptermodality{4}

\begin{chapteropening}
\textbf{By the end of this chapter:} \emph{I can say what I think, understand, read and write in Russian, name the finished partner each of those verbs travels with, and read the letter \ru{ш} that has been hiding in words I already say.}

Say all four verbs of the chapter in the I form — \ru{я} \ru{думаю}, \ru{я} \ru{понимаю}, \ru{я} \ru{читаю}, \ru{я} \ru{пишу} — then name each one's finished partner, put the stress on the end of \ru{писать} where it belongs, and read the \ru{ш} you have been decoding since \ru{живёшь}.
\end{chapteropening}

\section[dúmat']{\ru{думать} — \textquotedblleft{}to think\textquotedblright{}}

\label{lesson:RU-C04-dumat}



\begin{quote}
Six verbs behind you. This one opens the largest fact about Russian verbs.
\end{quote}

\subsection*{You'll want to know first — \ru{думать}}
\begin{quote}
\textbf{\ru{думать}} — \emph{dúmat'} — \textbf{to think}
\end{quote}

No new letters, and the \textbf{-\ru{ть}} of \emph{\ru{знать}}. Stress on the first syllable, and it never moves: \textbf{DOO-mat'}.

\begin{quote}
\textbf{\ru{Я} \ru{думаю}.} — \emph{ya dúmayu} — \textbf{I think.}
\end{quote}

That \textbf{-\ru{ю}} is the one you put on \emph{\ru{знаю}}. And \textbf{\ru{ты} \ru{думаешь}} — an \textbf{-\ru{ешь}}, so \emph{\ru{думать}} files with \emph{\ru{знать}}, not \emph{\ru{говорить}}.

Three verbs of the head now: \textbf{\ru{я} \ru{знаю}} is a fact I hold, \textbf{\ru{я} \ru{вижу}} is what reaches my eyes, \textbf{\ru{я} \ru{думаю}} is what goes on behind them.

\begin{grammarlens}[title={Grammar Lens: every Russian verb travels with a partner}]
Russian verbs come in \textbf{pairs}. One member is the \textbf{doing} — under way, repeated, unfinished. The other is the \textbf{deed done}, seen whole.

\emph{\ru{Думать}} is the first kind; its partner \textbf{\ru{подумать}} is not \textquotedblleft{}to think\textquotedblright{} but \textquotedblleft{}to have a thought, and be done with it.\textquotedblright{}

You have met the instinct once. \textbf{\ru{Идти}} was going \emph{now, one way}, with a separate partner for going \emph{as a habit} — one small family of motion verbs. This is a \textbf{different} distinction, applied to nearly \textbf{every verb in the language}.

So one job per verb: learn its partner's \textbf{name}. What a pair does to a sentence is a later lesson.
\end{grammarlens}

\begin{cousinweb}
\textbf{\ru{дума}} is the Russian noun for \textbf{a thought} — and for a body that thinks: the State \textbf{\ru{Дума}} is the lower house of the Russian parliament, this same word.

Its own origin carries a hedge. The standard account makes \emph{\ru{дума}} an early borrowing from \textbf{Gothic} \emph{dōms}, \textquotedblleft{}a judgement\textquotedblright{} — behind English \textbf{doom} (a judgement before it was a fate) and \textbf{deem}. If so, \emph{\ru{думать}} and \emph{doom} are relatives \textbf{by contact}, not by descent; a minority prefer an inherited root.

Latin's \emph{cogito, ergo sum} is built on the same idea, and its other pieces are yours: the \textbf{-o} of \emph{cogito} is Latin \textbf{ego}, the pronoun Russian keeps as \textbf{\ru{я}}; and \emph{fuī}, \textquotedblleft{}I was,\textquotedblright{} is \textbf{\ru{быть}}'s root, English \textbf{be}.
\end{cousinweb}

\subsection*{Your turn}
Say these aloud:

\begin{itemize}
\raggedright
  \item \textquotedblleft{}\ru{думать}\textquotedblright{} — \emph{DOO-mat'}, stress on the front
  \item the three heads — \textquotedblleft{}\ru{я} \ru{знаю} … \ru{я} \ru{вижу} … \ru{я} \ru{думаю}\textquotedblright{}
  \item the pair — \textquotedblleft{}\ru{думать} … \ru{подумать}\textquotedblright{}, the doing and the done
\end{itemize}

\begin{tcolorbox}[breakable,colback=teal!4,colframe=teal!35!black,title={Before you move on}]
Say \textquotedblleft{}I think.\textquotedblright{} (\textbf{\ru{Я} \ru{думаю}}.) Which family does \emph{\ru{думать}} join, \emph{\ru{знаешь}} or \emph{\ru{говоришь}}? (\textbf{\ru{Знаешь}}.) What is \emph{\ru{подумать}} beside \emph{\ru{думать}}, and how many verbs have such a partner? (The \textbf{finished} member; nearly \textbf{all}.) Which English word does the standard account attach \emph{\ru{дума}} to? (\textbf{Doom}, by borrowing.)
\end{tcolorbox}
\section[ponimát']{\ru{понимать} — \textquotedblleft{}to understand\textquotedblright{}}

\label{lesson:RU-C04-ponimat}



\begin{quote}
The most useful sentence a traveller owns is short, and after this lesson it is yours: \emph{I don't understand.}
\end{quote}

\subsection*{You'll want to know first — \ru{понимать}}
\begin{quote}
\textbf{\ru{понимать}} — \emph{ponimát'} — \textbf{to understand}
\end{quote}

No new letters. Stress on the last syllable, and both \textbf{\ru{о}}s lean towards \emph{a} away from it: \textbf{pa-ni-MAT'}.

\begin{quote}
\textbf{\ru{Я} \ru{понимаю}.} — \emph{ya panimáyu} — \textbf{I understand.}
\end{quote}

The same \textbf{-\ru{ю}} as \emph{\ru{думаю}}. And now the sentence worth more than the rest:

\begin{quote}
\textbf{\ru{Я} \ru{не} \ru{понимаю}.} — \emph{ya ne panimáyu} — \textbf{I don't understand.}
\end{quote}

\textbf{\ru{Не}} does the whole job, exactly as in \emph{\ru{я} \ru{не} \ru{знаю}}: one word in front of the verb, no helper borrowed from anywhere. Say those syllables to a Russian speaker and the conversation slows down for you.

\begin{grammarlens}[title={Grammar Lens: the partner, and a warning about its shape}]
\emph{\ru{Думать}} took \textbf{\ru{подумать}}: a prefix on the front, nothing else disturbed. Do not expect that every time.

\textbf{\ru{Понимать}}'s partner is \textbf{\ru{понять}} — same root, but the middle of the word has been squeezed out, \emph{-\ru{ним}-} against \emph{-\ru{ня}-}. Russian builds a pair by prefix, by a reshaped stem, or by both, and the shape is learned with the verb.

The job is unchanged: \textbf{\ru{понимать}, \ru{понять}}, said together, one entry.
\end{grammarlens}

\begin{cousinweb}
\textbf{\ru{понимать}} is \textbf{\ru{по}-} plus a root meaning simply \textbf{to take}: Indo-European *\emph{h\textsubscript{1}em-}, which came down the Latin branch as \emph{emere}, \textquotedblleft{}to take, and so to buy.\textquotedblright{} English took the family whole — \textbf{exempt} (taken out), \textbf{redeem} (bought back), \textbf{example} (taken as a specimen), \textbf{premium}, \textbf{prompt}, \textbf{assume}, \textbf{consume}.

So \emph{\ru{понимать}} literally says \textbf{take hold of}. English says the same and has stopped noticing: \textbf{comprehend} is Latin \emph{com-prehendere}, \textquotedblleft{}to seize together\textquotedblright{}; \textbf{grasp} is what a hand does; German \textbf{begreifen} is \emph{greifen}, \textquotedblleft{}to grip.\textquotedblright{}

Set that beside the two pictures Russian already gave you. \textbf{\ru{Знать}} was *\emph{ǵneh\textsubscript{3}-}, English \textbf{know} — knowledge you simply \emph{have}. \textbf{\ru{Видеть}} was *\emph{weid-}, which meant \emph{see} and \emph{know} at once and gave English \textbf{wit} and \textbf{wise} — knowledge you \emph{saw}. Now understanding as something you \emph{caught}. And unlike \emph{\ru{думать}}, which arrived by Gothic borrowing, this one is inherited straight down.
\end{cousinweb}

\subsection*{Your turn}
Say these aloud:

\begin{itemize}
\raggedright
  \item \textquotedblleft{}\ru{понимать}\textquotedblright{} — \emph{pa-ni-MAT'}, both \emph{o}s leaning to \emph{a}
  \item \textquotedblleft{}\ru{Я} \ru{понимаю}\textquotedblright{} then \textquotedblleft{}\ru{Я} \ru{не} \ru{понимаю}\textquotedblright{} — one word added, nothing else
  \item the pair — \textquotedblleft{}\ru{понимать} … \ru{понять}\textquotedblright{}, the squeezed middle
  \item the picture — \textquotedblleft{}take hold\textquotedblright{} … \emph{comprehend, grasp}
\end{itemize}

\begin{tcolorbox}[breakable,colback=teal!4,colframe=teal!35!black,title={Before you move on}]
Say \textquotedblleft{}I don't understand.\textquotedblright{} (\textbf{\ru{Я} \ru{не} \ru{понимаю}}.) Name \emph{\ru{понимать}}'s partner and say how its shape differs from \emph{\ru{подумать}}'s. (\textbf{\ru{Понять}} — the stem is reshaped, not just prefixed.) What does the root of \emph{\ru{понимать}} mean? (\textbf{To take} — as \emph{comprehend} and \emph{grasp} also do.) And which of your knowing-verbs is a borrowing rather than an inheritance? (\textbf{\ru{Думать}}, from Gothic.)
\end{tcolorbox}
\section[chitát']{\ru{читать} — \textquotedblleft{}to read\textquotedblright{}}

\label{lesson:RU-C04-chitat}



\begin{quote}
English says \emph{I read} and \emph{I am reading} and means two things. Russian cannot — and the reason is a verb it refuses to say.
\end{quote}

\subsection*{You'll want to know first — \ru{читать}}
\begin{quote}
\textbf{\ru{читать}} — \emph{chitát'} — \textbf{to read}
\end{quote}

No new letters: \textbf{\ru{ч}} you have been reading since \emph{\ru{очень}}, and it says the \emph{ch} of \emph{cheese}. Stress on the last syllable, \textbf{chi-TAT'}.

\begin{quote}
\textbf{\ru{Я} \ru{читаю}.} — \emph{ya chitáyu} — \textbf{I read.}
\end{quote}

The same \textbf{-\ru{ю}}, and \textbf{\ru{ты} \ru{читаешь}} puts it with \emph{\ru{знать}} and \emph{\ru{думать}}. Two verbs, one ordinary sentence:

\begin{quote}
\textbf{\ru{Я} \ru{читаю} \ru{и} \ru{понимаю}.} — \emph{ya chitáyu i panimáyu} — \textbf{I read and I understand.}
\end{quote}

\begin{grammarlens}[title={Grammar Lens: one present tense, two English sentences}]
\textbf{\ru{Я} \ru{читаю}} is \emph{I read} \textbf{and} \emph{I am reading}. Russian has one present tense and it covers both.

There is a reason, and you hold it. English builds \emph{I am reading} out of \textbf{be} plus \emph{-ing}; Russian's \textbf{\ru{быть}} has no present tense to build with — the missing verb of Chapter 3, the one \emph{\ru{Я} \ru{студент}} does without. With no auxiliary, Russian never grew the machinery.

The work English does with \emph{am reading} against \emph{read} is carried instead by the \textbf{pair}. \emph{\ru{Читать}} is the ongoing member; \textbf{\ru{прочитать}} is the finished one — to read a thing \textbf{through}, to the end.
\end{grammarlens}

\begin{cousinweb}
\textbf{\ru{читать}} rests on an old Slavic root meaning \textbf{to count, to reckon, to heed}, and Russian keeps the family in plain sight: \textbf{\ru{число}} is a \emph{number}, \textbf{\ru{честь}} is \emph{honour}, \textbf{\ru{почитать}} is \emph{to hold in esteem}. Counting, respecting and reading are one act — careful attention, item by item. Outside Slavic the root surfaces in Sanskrit \emph{cétati}, \textquotedblleft{}he perceives.\textquotedblright{}

And then the trail stops. There is \textbf{no secure English descendant}, and the two words that look promising are not it: \textbf{cheat} is a clipped \emph{escheat}, from Latin \emph{cadere}, \textquotedblleft{}to fall,\textquotedblright{} and \textbf{chit}, the little note, came into English from India. This book would rather hand you nothing than a cousin that is only a resemblance — the care it took with \emph{\ru{говорить}}, which is not \emph{govern}.
\end{cousinweb}

\subsection*{Your turn}
Say these aloud:

\begin{itemize}
\raggedright
  \item \textquotedblleft{}\ru{читать}\textquotedblright{} — \emph{chi-TAT'}, the \emph{ch} of \emph{cheese}
  \item \textquotedblleft{}\ru{Я} \ru{читаю} \ru{и} \ru{понимаю}\textquotedblright{} — and hear it as both \emph{read} and \emph{am reading}
  \item the pair — \textquotedblleft{}\ru{читать} … \ru{прочитать}\textquotedblright{}, the doing and the read-through
\end{itemize}

\begin{tcolorbox}[breakable,colback=teal!4,colframe=teal!35!black,title={Before you move on}]
Say \textquotedblleft{}I read and I understand.\textquotedblright{} (\textbf{\ru{Я} \ru{читаю} \ru{и} \ru{понимаю}}.) Give the two English sentences \emph{\ru{я} \ru{читаю}} covers, and say why. (\emph{I read} and \emph{I am reading}; \textbf{\ru{быть}} has no present tense to build the second.) Name the finished partner. (\textbf{\ru{Прочитать}}.) Which two Russian nouns share this root, and is \emph{\ru{читать}} related to English \emph{cheat}? (\textbf{\ru{Число}} and \textbf{\ru{честь}}; and \textbf{no} — as \emph{\ru{говорить}} was not related to \emph{govern}.)
\end{tcolorbox}
\section[pisát']{\ru{писать} — \textquotedblleft{}to write\textquotedblright{}}

\label{lesson:RU-C04-pisat}



\begin{quote}
A new letter, an English family you know on sight, and a stress mistake worth naming before you make it.
\end{quote}

\subsection*{You'll want to know first — \ru{писать}}
\begin{quote}
\textbf{\ru{писать}} — \emph{pisát'} — \textbf{to write}
\end{quote}

Stress on the \textbf{last} syllable, \textbf{pi-SAT'} — not a detail you may skip.

\begin{quote}
\textbf{\ru{Я} \ru{пишу}.} — \emph{ya pishú} — \textbf{I write}, and \textbf{I am writing}.
\end{quote}

Two things moved: the \textbf{\ru{с}} becomes \textbf{\ru{ш}} in the present, and the stress that sat on the ending for \emph{I} jumps back to the stem elsewhere — \textbf{\ru{пишу}} (\emph{pishú}), but \textbf{\ru{ты} \ru{пишешь}} (\emph{píshesh}). The pair runs as \emph{\ru{читать}}'s did: \textbf{\ru{писать}} ongoing, \textbf{\ru{написать}} finished.

\begin{quote}
\textbf{Careful.} \emph{\ru{Писать}} stressed at the end — \emph{pisát'} — is \textbf{to write}. \emph{\ru{Писать}} stressed at the front — \emph{písat'} — is a child's word for \textbf{to urinate}, whose \emph{I} form is \emph{písayu} beside \emph{pishú}. The two differ only in where your voice lifts, and beginners land on the wrong one. Stress the \textbf{end} and say \emph{pishú}.
\end{quote}

\subsection*{The letters in this word}
\emph{(Skim if you read Cyrillic.)} All of \textbf{\ru{писать}} is familiar; the new letter arrives in \textbf{\ru{пишу}}.

\begin{itemize}
\raggedright
  \item \textbf{\ru{ш}} = \textbf{\textquotedblleft{}sh\textquotedblright{}}, the sound of \emph{shoe}. Three uprights on a bar, like a comb with the teeth up.
\end{itemize}

You have been reading it without being told: \emph{\ru{живёшь}}, \emph{\ru{знаешь}}, \emph{\ru{говоришь}} and the past \emph{\ru{шёл}} all carry it. And note the kind of letter it is: \textbf{\ru{я}} lies to a Latin reader, looking like a mirrored \emph{R} while being a vowel; \textbf{\ru{ш}} resembles nothing Latin at all.

\begin{cousinweb}
\textbf{\ru{писать}} is Indo-European *\emph{pey\'{k}-}, \textquotedblleft{}\textbf{to cut, to scratch, to mark}\textquotedblright{} — before ink there was a point and a surface. The root became, in Latin, \emph{pingere}, \textquotedblleft{}to paint\textquotedblright{}, and English took the results: \textbf{paint}, \textbf{picture}, \textbf{pigment}. Russian's \textbf{\ru{письмо}} (a letter) and \textbf{\ru{писатель}} (a writer) sit on the verb.

That closes the chapter's thread: \textbf{\ru{думать}} came in from Gothic, \textbf{\ru{понимать}} is a root meaning \emph{take}, \textbf{\ru{читать}} a root meaning \emph{count} with no English descendant, and \textbf{\ru{писать}} a root meaning \emph{scratch} that handed English a wall of paint.
\end{cousinweb}

\subsection*{Your turn}
Say these aloud:

\begin{itemize}
\raggedright
  \item \textquotedblleft{}\ru{писать}\textquotedblright{} then \textquotedblleft{}\ru{я} \ru{пишу}\textquotedblright{} — stress on the end, and the \emph{\ru{ш}} arriving
  \item the four — \textquotedblleft{}\ru{я} \ru{думаю}, \ru{я} \ru{понимаю}, \ru{я} \ru{читаю}, \ru{я} \ru{пишу}\textquotedblright{}
\end{itemize}

\begin{tcolorbox}[breakable,colback=teal!4,colframe=teal!35!black,title={Before you move on}]
Say \textquotedblleft{}I write.\textquotedblright{} (\textbf{\ru{Я} \ru{пишу}}.) Where does the stress go, and what if it goes on the front? (On the \textbf{end}; \emph{písat'} is a child's word for urinating.) What sound is \textbf{\ru{ш}}? (\textbf{Sh} — \emph{\ru{живёшь}}, \emph{\ru{знаешь}}, \emph{\ru{говоришь}}, \emph{\ru{шёл}}.) Say the chapter's four verbs in the \emph{I} form, then their finished partners. (\textbf{\ru{Я} \ru{думаю}, \ru{я} \ru{понимаю}, \ru{я} \ru{читаю}, \ru{я} \ru{пишу}}; \textbf{\ru{подумать}, \ru{понять}, \ru{прочитать}, \ru{написать}}.) And what did \emph{\ru{писать}}'s root mean? (\textbf{To scratch} — hence \emph{paint}.)
\end{tcolorbox}
